\documentclass[11pt]{article}

\usepackage[margin=25mm]{geometry}

\usepackage{amsmath,amssymb,fontspec,unicode-math}

\IfFontExistsTF{DejaVu Serif}{\setmainfont{DejaVu Serif}}{\setmainfont{Latin Modern Roman}}

\IfFontExistsTF{Latin Modern Math}{\setmathfont{Latin Modern Math}}{\setmathfont{DejaVu Math TeX Gyre}}

\usepackage{xurl,hyperref}

\setlength{\parindent}{0pt}\setlength{\parskip}{6pt}

\begin{document}

\title{Sparse Image Synthesis via Joint Latent and RoI Flow}

\author{ELVIA preliminary AI analysis}\date{21 October 2026}\maketitle

Preliminary assessment: No Issue Found.\par

Independent scientific review remains pending. Author adoption has not been requested.

The augmented RoI loss changes the ideal conditional-mean drift, but the paper presents it as an empirical design choice and makes no exactness theorem for that loss. No substantive error was established in the inspected objective and schedule.

\section*{Review scope}

Rapid, claim-focused reading of the latest applicable author-paper and resource passages. Direct mathematical or static-source reasoning is stated below. No author program, retained verifier, model inference or benchmark was executed.

\section*{Problem 1: The extra absolute-error loss is an empirical objective choice}

Source: Proceedings Eqs.(7)–(13), PDF pp.5–6; Table 4, p.7; default loss weight, p.8

Assessment: not\_issue; impact: no\_issue\_found

\subsection*{Source claim (paraphrase)}

The method adds an \(L^1\) term to squared RoI-velocity regression. The paper does not claim a formal exactness theorem for this augmented loss.

\subsection*{Analysis}

The distinction is real but does not refute a stated guarantee: with the paper’s \(\beta=2\) schedule, independent Gaussian initial RoI and target values \(1/4,3/4\) of probabilities \(3/4,1/4\), the conditional initial velocities are \(1/2-2r,3/2-2r\). Squared regression has mean minimizer \(3/4-2r\); adding scalar absolute-error weight \(1/5\) gives the unique minimizer \(7/10-2r\). This shows that the usual conditional-mean derivation does not automatically transfer to the modified objective. It does not establish an incorrect final distribution or degraded image quality. Table 4’s reported \(L^1\) ablation supports its empirical use.

\subsection*{Finding and scope}

No substantive error is established within this focused check. Retain the mean-versus-augmented-objective distinction as scope, without demanding removal of an empirically chosen loss.

\section*{Limits}

Preliminary AI source/argument assessment; independent human scientific adjudication remains pending.\par

Mathematical counterexamples and source discrepancies do not remeasure published experiments; historical execution results were not rerun.\par

\section*{Source identity}

NeurIPS 2025 official proceedings PDF; current endpoint freshly retrieved and relevant current author revisions compared

\url{https://papers.nips.cc/paper_files/paper/2025/file/d6f7e8076b5da40e2c55582b3707419e-Paper-Conference.pdf}

Reviewed PDF SHA-256: \nolinkurl{7dbb044e98acf6128dd92f2491d7ed9a23ad32ac2151628c942fbb278fe76477}

The exact current PDF was freshly retrieved from the official endpoint. Its SHA-256 identifies the bytes inspected; the supplied historical LaTeX report remains unchanged outside this export. Paper authors were checked against the primary PDF header.

\end{document}
